% Source: https://artofproblemsolving.com/community/c3h1649341_product_of_divisors
% Stable identity: aops-topic-1649341-q01
\textbf{(a)} The map $d\mapsto\dfrac {n}{d}$ pairs the divisors of $n$ with themselves (it is its own inverse). Hence
\[
p(n)^2=\prod_{d\mid n}d\cdot\prod_{d\mid n}\dfrac {n}{d}=\prod_{d\mid n}\left(d\cdot\dfrac {n}{d}\right)=n^{\tau(n)} ,
\]
and since $p(n)>0$, $p(n)=n^{\dfrac{\tau(n)}{2}}$. (When $n$ is a perfect square, $\tau(n)$ is odd and $n^{\dfrac{\tau(n)}{2}}$ is a power of $\sqrt n$, an integer.)

\textbf{(b)} Dividing out the divisor $n$,
\[
\prod_{d\mid n,\ d<n}d=\dfrac{p(n)}{n}=\boxed{n^{\dfrac{\tau(n)}{2}-1}} .
\]
For $n=12$: the divisors are $1,2,3,4,6,12$, so $\tau(12)=6$ and $p(12)=1\cdot2\cdot3\cdot4\cdot6\cdot12=1728=12^3$; the proper divisors have product $144=12^{2}$.
